\chapter*{Предисловие}
\addcontentsline{toc}{chapter}{Предисловие}\markboth{Предисловие}{Предисловие}
\bookpage{3}

Повышенный интерес к теории $R$-функций (или, если говорить с более общих
позиций, теории $R$-отображений) объясняется как простотой и привлекательностью
ее начальных идей, так и широким диапазоном областей, в которых она находит
приложение. В этом отношении теория $R$-функций напоминает теорию групп: обе
эти теории базируются на чисто алгебраических концепциях, а их приложения
весьма многообразны. Для $R$-функций областями применения оказались:
аналитическая геометрия, математическое программирование, оптимальное
размещение геометрических объектов, распознавание образов, теория
устойчивости, химическая технология и в особенности математическая физика
с ее многочисленными задачами исследования, расчета и оптимизации
физико-механических полей (температурных, деформационных, силовых,
электромагнитных, магнитогидродинамических и~др.).

Цель монографии "--- изложить основные результаты, полученные почти за
двадцатилетний период в теории $R$-функций и ряде ее приложений. Особенно
интенсивно стала развиваться и применяться теория $R$-функций в последние
несколько лет. В связи с этим возникла необходимость осмысливания и обобщения
полученных результатов, однако вышедшие в свет монографии в основном были
посвящены отдельным прикладным направлениям и не отражали в должной мере
прогресс в развитии многих вопросов теории.

Монография состоит из введения и четырех глав. Гл.~1 вводит читателя в круг
представлений алгебры логики и исчисления предикатов (называемых также
дискретной или конечной математикой). Эти вопросы освещаются лишь в объеме,
необходимом для дальнейшего изложения. Более полно данные вопросы рассмотрены
в работе С.\,В.~Яблонского~[194]. В гл.~2 вводятся понятия $R$-отображения
($R$-функции), чертежа, уравнения чертежа, алгебраически полной системы и так
далее (термин <<чертеж>> используется вместо часто употребляемых терминов
<<линия>>, <<фигура>>, <<поверхность>>, <<тело>>, <<гиперповерхность>>,
<<геометрический объект>> и~т.\,п.). В этой же главе решается принципиально
важный вопрос о построении уравнений чертежей (решение обратной задачи
аналитической геометрии), рассматриваются приложения полученных результатов
в математическом программировании и оптимальном размещении геометрических
объектов. Новый в идейном плане материал в основном содержится в гл.~3 и~4.

Материал гл.~3, в которой используются результаты предыдущих глав, посвящен
развитию конструктивных средств математики и решению казавшейся ранее
неразрешимой проблемы построения полных координатных последовательностей,
удовлетворяющих различного типа краевым условиям для областей
\bookpage{4}%
практически произвольной формы. Эти результаты легли в основу новой
технологии программирования и ее реализации в виде программирующих систем
(генераторов программ) серии <<Поле>>, позволивших в десятки раз сократить
затраты на составление и отладку программ. О генераторах программ и решении
краевых задач в условиях их эксплуатации речь идет в гл.~4. Материал этой
главы подобран с таким расчетом, чтобы не только показать возможности
применения $R$-функций к решению задач расчета полей различного типа, но и
отметить некоторые специфические особенности реализации теории $R$-функций на
практике.

Многие параграфы сопровождаются упражнениями и задачами. Некоторые из них
носят проблемный характер, и автор не знает их решений. Ряд нерешенных
вопросов отмечен в основном тексте. Большое число примеров, приведенных
в работе, облегчает освоение материала, делает его доступным для студентов
старших курсов инженерных и математических специальностей. (Примеры
и некоторые параграфы, которые при первом чтении можно опустить, набраны
петитом.) В списке литературы указаны лишь те источники, содержание которых
в какой-то мере связано с излагаемыми результатами. Обширная библиография
содержится в монографиях и статьях автора и его учеников И.\,В.~Гончарюка,
В.\,С.~Проценко, Ю.\,Г.~Стояна, В.\,А.~Рвачева, А.\,П.~Слесаренко,
Г.\,П.~Манько, Т.\,И.~Шейко, Н.\,С.~Синекопа, Ф.\,Ф.~Коваля и~др. [25--27, 49,
50, 67--75, 100, 101, 109--148, 154--156, 160--164, 167--172, 184--188].
Ссылки (на теоремы, примеры, упражнения и~т.\,д.) внутри параграфов (глав)
даются без указаний номеров параграфов (глав).
